% Book-style note WITHOUT QR (from slides / slides+transcript) — compile: scripts/build.sh note-noqr.tex
\documentclass[10pt,twoside]{article}
\usepackage[tiro]{edtech}
\geometry{left=13mm,right=13mm,top=17mm,bottom=17mm,headsep=5mm}
\etset{brand=EduBrand, book={HSC Physics 2nd Paper · Chapter-03},
       chapter={চল তড়িৎ}, logo=logo-placeholder.png}
\pagestyle{note}
\setlength{\columnseprule}{0.4pt}\setlength{\columnsep}{6mm}
\ctikzset{bipoles/length=0.9cm}
\begin{document}
\chapterbanner{তৃতীয় অধ্যায়}{চল তড়িৎ}{Current Electricity}{\etlogo[1.1cm]}

\sectionbar{\centering Board Questions Analysis}
\textbf{সৃজনশীল প্রশ্ন}\hfill{\scriptsize\color{etgraydark}* ডেমো ডেটা — আসল বিশ্লেষণে আপনার দেওয়া বোর্ড-তথ্য বসবে}\par\smallskip
{\renewcommand\arraystretch{1.3}\small
\noindent\begin{tabularx}{\linewidth}{|>{\raggedright\arraybackslash}p{0.3\linewidth}|c|c|c|X|c|}\hline
\rowcolor{etgray} & & \multicolumn{2}{c|}{\textbf{যতবার এসেছে}} & & \\\cline{3-4}
\rowcolor{etgray}\multirow{-2}{*}{\textbf{টাইপের নাম}} & \multirow{-2}{*}{\textbf{গুরুত্ব}} & ‘গ’ & ‘ঘ’ & \multirow{-2}{*}{\textbf{CQ প্রশ্নসমূহ যে বোর্ডে যে বছর এসেছে}} & \multirow{-2}{*}{\textbf{পৃষ্ঠা}}\\\hline
{\latinserif\bfseries Type-01}: তুল্য রোধ ও কার্শফের সূত্র & \stars{5} & ১৮ & ১৫ & ঢা. বো. ২৫; রা. বো. ২৫; চ. বো. ২৪; য. বো. ২৩; সি. বো. ২৩; দি. বো. ২২ & ২\\\hline
{\latinserif\bfseries Type-02}: হুইটস্টোন ব্রিজ ও মিটার ব্রিজ & \stars{4} & ১১ & ৯ & কু. বো. ২৫; ব. বো. ২৪; ম. বো. ২৪; ঢা. বো. ২৩; রা. বো. ২২ & ২\\\hline
{\latinserif\bfseries Type-03}: তড়িৎ ক্ষমতা ও জুলের সূত্র & \stars{4} & ৯ & ৭ & সি. বো. ২৫; চ. বো. ২৩; য. বো. ২২; সকল বোর্ড ১৮ & ৩\\\hline
\end{tabularx}}

\slantedtitle{গুরুত্বপূর্ণ সূত্রাবলি ও বিশ্লেষণ}
{\renewcommand\arraystretch{1.15}
\noindent\begin{tabularx}{\linewidth}{|>{\raggedright}p{0.5\linewidth}|>{\centering\arraybackslash}X|}\hline
\rowcolor{etgray}\multicolumn{1}{|c|}{\headfont\bfseries সূত্রাবলি} & {\headfont\bfseries বিশ্লেষণ}\\\hline
$\blacksquare$ \textbf{ওহমের সূত্র:}\ $V=IR$\\[2pt]
 \hspace*{1em}$V=$ বিভব পার্থক্য,\ $I=$ তড়িৎ প্রবাহ,\ $R=$ রোধ\\[2pt]
 $\blacksquare$ \textbf{আপেক্ষিক রোধ:}\ $R=\rho\dfrac{L}{A}$
&\begin{circuitikz}[scale=0.65,baseline={([yshift=-2pt]current bounding box.north)},font=\footnotesize]
  \draw (0,0) to[R,l=$R$,i>^=$I$] (3,0); \draw[{Latex}-{Latex}](0,-0.6)--(3,-0.6) node[midway,below]{$V$};
 \end{circuitikz}\tabularnewline\hline
$\blacksquare$ \textbf{শ্রেণি সমবায়:}\ $R_s=R_1+R_2+R_3$\\[2pt]
 \hspace*{1em}প্রতিটি রোধে প্রবাহ সমান; বিভব ভাগ হয়।\\[3pt]
 $\blacksquare$ \textbf{সমান্তরাল সমবায়:}\ $\dfrac{1}{R_p}=\dfrac{1}{R_1}+\dfrac{1}{R_2}+\dfrac{1}{R_3}$\\[2pt]
 \hspace*{1em}প্রতিটি রোধে বিভব সমান; প্রবাহ ভাগ হয়।
&\begin{circuitikz}[scale=0.55,baseline={([yshift=-2pt]current bounding box.north)},font=\scriptsize]
  \draw (0,3.0) to[R,l=$R_1$] (1.7,3.0) to[R,l=$R_2$] (3.4,3.0) to[R,l=$R_3$] (5.1,3.0);
  \node[font=\scriptsize] at (2.55,2.35){শ্রেণি};
  \draw (0.4,1.5) to[R,l=$R_1$] (2.6,1.5) (0.4,0) to[R,l=$R_2$] (2.6,0) (0.4,-1.5) to[R,l=$R_3$] (2.6,-1.5);
  \draw (0.4,1.5)--(0.4,-1.5) (2.6,1.5)--(2.6,-1.5) (-0.3,0)--(0.4,0) (2.6,0)--(3.3,0);
  \node[font=\scriptsize,right] at (3.3,0){সমান্তরাল};
 \end{circuitikz}\tabularnewline\hline
$\blacksquare$ \textbf{কার্শফের সূত্র:}\\[2pt]
 $\bullet$ প্রথম সূত্র (জংশন):\ $\sum I=0$\\[2pt]
 $\bullet$ দ্বিতীয় সূত্র (বদ্ধ লুপ):\ $\sum E=\sum IR$
&\begin{circuitikz}[scale=0.55,baseline={([yshift=-2pt]current bounding box.north)},font=\scriptsize]
  \draw (0,0) to[battery1,l=$E_1$] (0,2.2) to[R,l=$R_1$] (2.4,2.2) to[R,l=$R_2$] (2.4,0) -- (0,0);
  \draw (2.4,2.2) to[R,l=$R_3$] (4.8,2.2) to[battery1,l=$E_2$] (4.8,0) -- (2.4,0);
  \draw[-{Latex[length=1.6mm]},etblue] (1.55,1.45) arc (90:-200:0.42);
  \node[circle,fill,inner sep=1.2pt] at (2.4,2.2){};\node[circle,fill,inner sep=1.2pt] at (2.4,0){};
 \end{circuitikz}\tabularnewline\hline
$\blacksquare$ \textbf{হুইটস্টোন ব্রিজ:}\\[2pt]
 সাম্যাবস্থায় (গ্যালভানোমিটারে প্রবাহ শূন্য)\\[2pt]
 \hspace*{2em}$\dfrac{P}{Q}=\dfrac{R}{S}$\\[2pt]
 $\blacksquare$ \textbf{কোষের প্রবাহ:}\ $I=\dfrac{E}{R+r}$\quad ($r=$ অভ্যন্তরীণ রোধ)
&\begin{circuitikz}[scale=0.55,baseline={([yshift=-2pt]current bounding box.north)},font=\scriptsize]
  \draw (0,0) to[R,l=$P$] (1.8,1.8) to[R,l=$Q$] (3.6,0);
  \draw (0,0) to[R,l_=$R$] (1.8,-1.8) to[R,l_=$S$] (3.6,0);
  \draw (1.8,1.8) to[rmeter,t=G] (1.8,-1.8);
  \draw (0,0) -- (-0.6,0) -- (-0.6,-2.6) to[battery1,l_=$E$] (4.2,-2.6) -- (4.2,0) -- (3.6,0);
 \end{circuitikz}\tabularnewline\hline
\end{tabularx}}

\begin{remember}[একক রূপান্তর]
\noindent$\blacklozenge$\ $\SI{1}{kWh}=\SI{3.6e6}{J}$ (১ ইউনিট)\qquad
$\blacklozenge$\ $\SI{1}{\micro\ohm}=10^{-6}\,\Omega$\qquad
$\blacklozenge$\ $P=VI=I^2R=\dfrac{V^2}{R}$\qquad
$\blacklozenge$\ $H=I^2Rt$ (জুল)
\end{remember}

\newpage
\begin{multicols}{2}
\typehead{01}{তুল্য রোধ ও কার্শফের সূত্র}
\begin{shadowbox}
$\blacksquare$ \textbf{কৌশল:} জটিল বর্তনীতে আগে \textbf{সমান্তরাল অংশ} মিলিয়ে একটি রোধ বানাও, তারপর শ্রেণিতে যোগ কর। কোষের অভ্যন্তরীণ রোধ $r$ থাকলে সেটিও শ্রেণিতে যোগ হবে—
$$I=\frac{E}{R_{\text{তুল্য}}+r}$$
টার্মিনাল বিভব:\ $V=E-Ir$
\end{shadowbox}
\medskip
\prashna{১} নিচের বর্তনীটি লক্ষ কর: \boardtag{উদ্দীপক}\par
\begin{center}\begin{circuitikz}[scale=0.75,font=\scriptsize]
 \draw (0,0) to[battery1,l={$\SI{12}{V}$, $r=\SI{1}{\ohm}$}] (0,2.4) to[R,l=$\SI{4}{\ohm}$] (2.6,2.4) -- (3.2,2.4);
 \draw (3.2,2.4) -- (3.2,0)  (3.2,2.4) -- (4.6,2.4) to[R,l=$\SI{6}{\ohm}$] (4.6,0) -- (3.2,0);
 \draw (3.2,2.4) to[R,l_=$\SI{3}{\ohm}$,*-*] (3.2,0);
 \draw (3.2,0) -- (0,0);
\end{circuitikz}\end{center}
\textbf{(গ)} বর্তনীর মূল প্রবাহ নির্ণয় কর।\\
\textbf{(ঘ)} \SI{6}{\ohm} রোধটি খুলে নিলে \SI{4}{\ohm} রোধের দুই প্রান্তের বিভব পার্থক্য বাড়বে নাকি কমবে? গাণিতিকভাবে বিশ্লেষণ কর।\par\smallskip
\textbf{সমাধান:}\par
\partlabel{গ}
\solgiven[0.38]{
$R_p=\dfrac{6\times3}{6+3}=\SI{2}{\ohm}$\\[2pt]
$R_{\text{তুল্য}}=4+2=\SI{6}{\ohm}$\\[2pt]
$I=\dfrac{E}{R_{\text{তুল্য}}+r}=\dfrac{12}{6+1}$\\[2pt]
$=\SI{1.71}{A}$\ans}
{\textbf{দেওয়া আছে,}\\ $E=\SI{12}{V}$\\ $r=\SI{1}{\ohm}$\\ $R_1=\SI{4}{\ohm}$\\ $R_2=\SI{6}{\ohm}$, $R_3=\SI{3}{\ohm}$}
\partlabel{ঘ}\ প্রথমে \SI{4}{\ohm} রোধে বিভব পার্থক্য,\\
$V_1=IR_1=1.714\times4=\SI{6.86}{V}$\\[2pt]
\SI{6}{\ohm} খুলে নিলে, $R'_{\text{তুল্য}}=4+3=\SI{7}{\ohm}$\\[2pt]
$I'=\dfrac{12}{7+1}=\SI{1.5}{A}$,\quad $V_1'=1.5\times4=\SI{6}{V}$\\[2pt]
$\therefore$ $V_1'<V_1$; অর্থাৎ বিভব পার্থক্য \textbf{কমবে}।\ans

\columnbreak
\typehead{02}{হুইটস্টোন ব্রিজ}
\begin{conceptbox}
\textbf{কনসেপ্ট:} চারটি রোধ $P,Q,R,S$ দিয়ে তৈরি ব্রিজের মধ্যবর্তী শাখায় গ্যালভানোমিটার থাকে। $\dfrac{P}{Q}=\dfrac{R}{S}$ হলে $B$ ও $D$ বিন্দুর বিভব সমান হয়, তাই গ্যালভানোমিটারে \textbf{কোনো প্রবাহ থাকে না}।
\begin{roundbox}[left=3pt,top=2pt,bottom=2pt]
$\bullet$ সাম্যাবস্থায় কোষ ও গ্যালভানোমিটারের অবস্থান অদলবদল করলেও সাম্য নষ্ট হয় না।\\
$\bullet$ অজানা রোধ নির্ণয়ে মিটার ব্রিজ এই নীতিতেই কাজ করে।
\end{roundbox}
\end{conceptbox}
\medskip
\prashna{২} একটি হুইটস্টোন ব্রিজে $P=\SI{10}{\ohm}$, $Q=\SI{20}{\ohm}$, $R=\SI{15}{\ohm}$ এবং $S=\SI{25}{\ohm}$।\par
\textbf{(গ)} ব্রিজটি সাম্যাবস্থায় আছে কি না যাচাই কর।\\
\textbf{(ঘ)} ব্রিজটিকে সাম্যাবস্থায় আনতে $S$ এর সাথে কত মানের রোধ কীভাবে যুক্ত করতে হবে? বিশ্লেষণ কর।\par\smallskip
\textbf{সমাধান:}\par
\partlabel{গ}
\solgiven[0.36]{
$\dfrac{P}{Q}=\dfrac{10}{20}=0.5$\\[2pt]
$\dfrac{R}{S}=\dfrac{15}{25}=0.6$\\[2pt]
$\dfrac{P}{Q}\neq\dfrac{R}{S}$; সাম্যাবস্থায় নেই।\ans}
{\textbf{দেওয়া আছে,}\\ $P=\SI{10}{\ohm}$, $Q=\SI{20}{\ohm}$\\ $R=\SI{15}{\ohm}$, $S=\SI{25}{\ohm}$}
\partlabel{ঘ}\ সাম্যাবস্থার জন্য প্রয়োজনীয় রোধ,\\
$S'=\dfrac{QR}{P}=\dfrac{20\times15}{10}=\SI{30}{\ohm}$\\[2pt]
যেহেতু $S'>S$, তাই $S$ এর সাথে \textbf{শ্রেণিতে} রোধ যোগ করতে হবে—\\
$S'-S=30-25=\SI{5}{\ohm}$\\[2pt]
$\therefore$ $S$ এর সাথে শ্রেণিতে \SI{5}{\ohm} রোধ যুক্ত করতে হবে।\ans
\end{multicols}
\end{document}
